% Discussion 4.2 --- counterintuitive environmental and thermal signals.
\begin{table}[htbp]
\centering
\caption{Counterintuitive environmental and thermal signals.}
\label{tab:insights-counterintuitive}
\footnotesize
\renewcommand{\arraystretch}{1.3}
\begin{tabularx}{\textwidth}{@{}>{\raggedright\arraybackslash\bfseries}p{0.30\textwidth} >{\raggedright\arraybackslash}X@{}}
\toprule
\normalfont\textbf{Key insight} & \textbf{Detail} \\
\midrule
Thermal preference is the weakest spatial signal & Its association with space type (Cram\'er's V $\approx 0.12$) sits far below activity or noise. \\
Cooling demand separates spaces only modestly & Outdoor demand ($\approx 53\%$) exceeds home ($\approx 43\%$) by less than an air-conditioning-driven intuition would suggest. \\
Offices can over-cool & Wanting \emph{warmer}, rare overall, concentrates in offices (lift $\approx 1.6$), a signature of air-conditioning overshoot. \\
The kind of noise matters more than the amount & The source of noise fingerprints a space more sharply than how loud it is. \\
Lift exposes hidden distinctiveness & \emph{A lot} of noise is $\approx 2.6\times$ over-represented in transportation relative to the whole corpus, a contrast raw proportions understate. \\
\bottomrule
\end{tabularx}
\end{table}
